% Preámbulo compartido para los modelos del segundo parcial de ACN.
% Replica el formato del parcial real 2P-2S-2025.
\usepackage[utf8]{inputenc}
\usepackage[T1]{fontenc}
\usepackage[spanish]{babel}
\usepackage{amsmath,amssymb}
\usepackage[margin=2.6cm]{geometry}
\usepackage{enumitem}
\usepackage{array}

\pagestyle{plain}
\setlength{\parindent}{0pt}
\setlength{\parskip}{5pt}

% Listas de incisos: a), b), c), ...
\setlist[enumerate,1]{label=\alph*),leftmargin=2.2em,itemsep=7pt,topsep=7pt}
% Listas numeradas 1., 2., 3., ...
\newlist{puntos}{enumerate}{1}
\setlist[puntos]{label=\arabic*.,leftmargin=2.2em,itemsep=7pt,topsep=7pt}
% Listas de condiciones del enunciado
\setlist[itemize,1]{label=$\bullet$,leftmargin=2em,itemsep=5pt,topsep=6pt}

% Encabezado del parcial
\newcommand{\encabezado}[1]{%
  \begin{center}
    \textbf{\large Aplicaciones Computacionales en Negocios}\\[3pt]
    \textbf{\large Segundo semestre de 2025 --- Segundo parcial}\\[3pt]
    \textbf{\normalsize Modelo de práctica #1}
  \end{center}
  \vspace{4pt}
  \hrule height 0.8pt
  \vspace{14pt}
  Nombre:\par\vspace{12pt}
  Legajo:\par\vspace{12pt}
  DNI:\par\vspace{12pt}
  ¿Te recibís con este parcial?:
  \vspace{10pt}
  \hrule height 0.8pt
  \vspace{16pt}
}

% Título de ejercicio
\newcommand{\ejercicio}[1]{\par\medskip\noindent\textbf{Ejercicio #1.}\ }

% Atajos de notación usados por la cátedra
\newcommand{\R}{\mathbb{R}}
\newcommand{\Z}{\mathbb{Z}}
\newcommand{\dem}{\mathrm{dem}}
\newcommand{\capa}{\mathrm{cap}}
\newcommand{\costo}{\mathrm{costo}}
